\documentclass[11pt,a4paper]{scrartcl}
\usepackage{iutparakou}
\begin{document}

\fichetitre{Jour 5 — Intégration \& Myhill--Nerode}{}
{Modules : \texttt{pipeline.py}, \texttt{myhill\_nerode.py} \quad$\bullet$\quad Barème : 3 pts (+1)}

\section*{Contexte}
Deux visages. \textbf{Assembler} : les étages, jusqu'ici isolés, \textbf{coopèrent} dans
un \texttt{pipeline}. \textbf{Boucler la boucle théorique} : on revient sur la question
du jour~1 — « combien d'états faut-il vraiment ? » — avec l'outil qui y répond
exactement, la \textbf{congruence de Myhill--Nerode}.

\section*{Notions nouvelles du jour}

\subsection*{1. Intégration : deux chaînes, un même cœur}
Le \texttt{pipeline} \textbf{appelle les apps}, qui \textbf{instancient} \texttt{formlang}.
Deux chaînes de traitement :

\begin{center}
\begin{tikzpicture}[etape/.style={draw=iutblue,very thick,rounded corners=3pt,
    fill=iutblue!6,minimum height=7mm,inner xsep=5pt,font=\footnotesize\bfseries,text=iutblue,align=center},
    fl/.style={-{Stealth[length=2mm]},very thick,iutorange!85!black}, node distance=8mm]
  % chaîne mot
  \node[etape] (fst) {FST\\normalise};
  \node[etape,right=of fst] (dfa) {AFD\\facteur \texttt{or}};
  \node[etape,right=of dfa] (pda) {PDA\\délimiteurs};
  \node[left=6mm of fst,font=\scriptsize\itshape,text=iutgrey] {\texttt{analyze\_word}};
  \draw[fl] (fst)--(dfa); \draw[fl] (dfa)--(pda);
  % chaîne morpho
  \node[etape,below=9mm of fst] (disc) {découverte\\affixes};
  \node[etape,right=of disc] (seg) {segmentation};
  \node[etape,right=of seg] (buta) {BUTA\\classify};
  \node[left=6mm of disc,font=\scriptsize\itshape,text=iutgrey] {\texttt{analyze\_morpho}};
  \draw[fl] (disc)--(seg); \draw[fl] (seg)--(buta);
\end{tikzpicture}
\end{center}
\noindent On évalue l'\textbf{architecture} (chaque étage instancie le cœur), pas un
nouvel algorithme.

\subsection*{2. La congruence de Myhill--Nerode}
$u\approx_{L}v$ si, \textbf{pour tout} suffixe $w$, $uw\in L\Leftrightarrow vw\in L$ :
une fois $u$ ou $v$ lu, l'\emph{avenir} (ce qu'on pourra encore accepter) est identique.
Équivalence \textbf{invariante à droite} — d'où l'on fait les \textbf{états} d'un automate.

\begin{notion}[Le théorème — le fil rouge se referme]
$L$ est régulier \textbf{ssi} $\approx_{L}$ a un nombre \textbf{fini} de classes, et ce
nombre (l'\textbf{indice}) est \textbf{exactement} le nombre d'états de l'AFD
\textbf{minimal}. Minimiser (jour~1), c'était calculer ces classes.
\end{notion}

\subsection*{3. Le calcul en pratique : signatures sur suffixes témoins}
On approxime « pour tout suffixe » par un \textbf{jeu fini} de suffixes distinctifs
$S=\{\varepsilon,\,r,\,or\}$. La \textbf{signature} d'un mot est
$\big(\mathbf{1}[w s\in L_1]\big)_{s\in S}$. Deux mots de même signature sont dans la même
classe :

\begin{exemple}[Trois signatures = trois classes]
\begin{center}\small
\begin{tabular}{@{}lccc>{\bfseries}c@{}}
\toprule
mot $w$ & $w\varepsilon$ & $wr$ & $wor$ & classe\\\midrule
\texttt{""}  & 0 & 0 & 1 & A\\
\texttt{a}   & 0 & 0 & 1 & A\\
\texttt{o}   & 0 & 1 & 1 & B\\
\texttt{ao}  & 0 & 1 & 1 & B\\
\texttt{or}  & 1 & 1 & 1 & C\\
\texttt{aor} & 1 & 1 & 1 & C\\
\bottomrule
\end{tabular}
\end{center}
(\texttt{1} = « contient \texttt{or} ».) Signatures $(0,0,1),(0,1,1),(1,1,1)$ : exactement
\textbf{3 classes} $\Rightarrow$ 3 états de l'AFD minimal. \texttt{nerode\_classes} fait
ce regroupement (test \texttt{test\_nerode\_trois\_classes}).
\end{exemple}

\begin{center}
\begin{tikzpicture}[classe/.style={draw=iutblue,very thick,rounded corners=4pt,
    fill=iutblue!5,minimum width=33mm,minimum height=14mm,align=center,font=\small}]
  \node[classe] (A) {\textbf{$C_A$}\\ \texttt{""},\ \texttt{a}\\[-1pt]{\footnotesize\bfseries\color{iutorange!85!black}état A}};
  \node[classe,right=8mm of A] (B) {\textbf{$C_B$}\\ \texttt{o},\ \texttt{ao}\\[-1pt]{\footnotesize\bfseries\color{iutorange!85!black}état B ($o$ armé)}};
  \node[classe,right=8mm of B] (C) {\textbf{$C_C$}\\ \texttt{or},\ \texttt{aor}\\[-1pt]{\footnotesize\bfseries\color{iutorange!85!black}état C (absorbant)}};
  \draw[-{Stealth[length=2mm]},thick,iutblue!70] (A) -- node[above,font=\scriptsize]{$o$} (B);
  \draw[-{Stealth[length=2mm]},thick,iutblue!70] (B) -- node[above,font=\scriptsize]{$r$} (C);
\end{tikzpicture}
\end{center}

\subsection*{4. Ouverture : Myhill--Nerode pour les arbres (bonus)}
Deux sous-arbres sont équivalents s'ils donnent le même verdict \textbf{dans tout
contexte} (arbre « à trou ») : fondement de la minimisation d'un BUTA, et lien direct
avec le hash-consing du jour~3 (partager, c'est identifier des sous-structures équivalentes).

\section*{A. Théorie \normalfont(à rédiger) — 1 pt}
\begin{description}[leftmargin=1.6em,style=nextline]
  \item[Q5.1 (0,5)] Énoncer Myhill--Nerode ; relier le nombre de classes de $L_1$ à
    l'AFD minimal.
  \item[Q5.2 (0,5)] Critère « indistinguables par tous les suffixes témoins » et lien
    avec la fusion d'états.
\end{description}

\section*{B. Pratique — 2 pts}
\begin{description}[leftmargin=1.6em,style=nextline]
  \item[E5.1 (1)] \texttt{analyze\_word} (FST$\to$AFD$\to$PDA) et \texttt{analyze\_morpho}
    (découverte$\to$segmentation$\to$BUTA). \emph{Tests pipeline.}
  \item[E5.3 (1)] \texttt{nerode\_classes} : regrouper par signature de suffixes.
    \emph{Test : \texttt{test\_nerode\_trois\_classes}.}
\end{description}

\begin{exemple}[Sortie de référence]
\ttfamily \$ pytest -q\\ 28 passed\\[2pt]
\$ python pipeline.py --word 4or\\
\hspace*{2em}normalisé(FST) : aor\ \ facteur\_or(AFD) : True\ \ délimiteurs\_ok(PDA) : True
\end{exemple}

\section*{Mini-barème}
\begin{center}\small
\begin{tabular}{@{}llr@{}}
\toprule Item & Détail & Pts\\\midrule
Q5.1–Q5.2 & Myhill--Nerode + lien minimisation & 1\\
E5.1 & pipeline (mot + morpho) & 1\\
E5.3 & classes de Nerode & 1\\
\textbf{Bonus} & mesure hash-consing sur 10\,000 mots & +1\\
\bottomrule
\end{tabular}
\end{center}

\begin{gate}
Le pipeline \textbf{appelle les apps} (qui instancient le cœur) ; il ne reprogramme
aucun automate.
\end{gate}

\subsection*{Références}
Hopcroft, Motwani, Ullman — Myhill--Nerode, minimisation. \textbullet\ Comon \emph{et
al.}, \emph{TATA} (2007) — congruence pour les arbres (bonus).

\end{document}
